\par\bigskip\Needspace{5\baselineskip}
\begin{center}\bfseries 4.\end{center}
\noindent\hspace*{\parindent}Der Fortgang von den Kongruenzen ersten zu denjenigen höheren Grades führt von selbst zur Theorie der Potenzreste d. i. zur Untersuchung der Reste, welche die Potenzen der verschiedenen Zahlen lassen, wenn sie durch einen gegebenen Modul geteilt werden, sowie der Bedingungen, unter denen eine gegebene Zahl ein solcher Rest einer Potenz von vorgeschriebenem Grade sein kann, u. a. mehr. Besonders interessiert hierbei der Fall, wo der Modul eine ungerade Primzahl \(p\) ist; \textsc{Gauss} behandelt daher \textls[100]{im dritten Abschnitte} diesen vorzugsweise. Für jede durch \(p\) nicht teilbare Zahl \(a\) gibt es eine kleinste Potenz \(a^d\), welche den Rest 1 lässt; dieser Exponent \(d\), zu welchem »\(a\ (\mathrm{mod.}\ p)\) gehört«, ist stets ein Teiler von \(p-1\), und aus diesem Umstande folgt sogleich der berühmte, schon von \textsc{Fermat}\footnote{\textsc{Fermatii} Opera mathem., Tolosae 1679, S. 163, Oeuvres II, S. 209.} ausgesprochene Satz, dass \(a^{p-1}\) stets kongruent 1 ist \((\mathrm{mod.}\ p)\). Der Beweis dieser Sätze durch die Methode der Exhaustion, wie \textsc{Gauss} ihn gibt, war durch eine im wesentlichen gleiche Betrachtung schon von \textsc{Euler}\begingroup\renewcommand{\thefootnote}{**)}\footnote{Novi Comment. Acad. Petropol. \textbf{7} (1758/59) 1761, S. 49, Opera Omnia, Ser. I, vol. 2, S. 493.}\endgroup\ erhalten, und schon früher\begingroup\renewcommand{\thefootnote}{***)}\footnote{Comment. Acad. Petropol. \textbf{8} (1736) 1741, S. 141, Opera Omnia, Ser. I, vol. 2, S. 33 und Novi Comment. Acad. Petropol. \textbf{1} (1747/8) 1750, S. 20, Opera Omnia, Ser. I, vol. 2, S. 62.}\endgroup\ hatte dieser den \textsc{Fermat}schen Satz mittels des binomischen Satzes begründet. \textsc{Gauss} fügte dem einen neuen Beweis mit Hilfe des polynomischen Satzes hinzu.

Wesentlich neu aber war die Bestimmung der Anzahl aller inkongruenten Zahlen, die zu einem gegebenen Teiler \(d\) von \(p-1\) als Exponenten gehören, aus welcher dann sofort die Tatsache der Existenz einer primitiven Wurzel \((\mathrm{mod.}\ p)\), d. i. einer Zahl \(g\), welche zum Exponenten \(p-1\) gehört, hervorging, eine Tatsache, welche nur irrtümlicherweise vordem \textsc{Euler} schon festgestellt zu haben gemeint hatte\begingroup\renewcommand{\thefootnote}{†)}\footnote{Novi Comment. Acad. Petropol. \textbf{18} (1773) 1774, S. 85.}\endgroup. Die Potenzen
\[
1,g,g^2,\ldots,g^{p-2}
\]
\clearpage\SourceLeaf{18}\ArithmeticHead{11}{R}{INHALT DER ABSCHNITTE I—IV DER D. A.}
\noindent einer solchen Wurzel haben die Eigenschaft, dass jede durch \(p\) nicht teilbare Zahl einer einzigen von ihnen \((\mathrm{mod.}\ p)\) kongruent ist. Dabei bestimmen die geraden Potenzen dieser Reihe die quadratischen Reste \((\mathrm{mod.}\ p)\), d. i. die Zahlen, welche \((\mathrm{mod.}\ p)\) einer Quadratzahl kongruent sein können. Aus dieser Bemerkung schloss \textsc{Euler}\footnote{Opuscula analytica I (1783), S. 242, 268, vergl. aber Novi Comment. Acad. Petropol. \textbf{1} (1747/8) 1750, S. 20 und \textbf{7} (1758/9) 1761, S. 49, Opera Omnia, Ser. I, vol. 2, S. 62 und 493, wo in den Theorematis 11 der ersten und 19 der zweiten Abhandlung das \textsc{Euler}sche Kriterium ohne Zuhilfenahme einer primitiven Wurzel vollkommen streng begründet wird.} — dem eben Gesagten zufolge aber ohne genügende Grundlage — sein bekanntes Kriterium für den quadratischen Charakter einer Zahl \(a\), nach welchem
\[
a^{\frac{p-1}{2}}\equiv\pm1\ (\mathrm{mod.}\ p)
\]
ist, jenachdem \(a\) ein quadratischer Rest \((\mathrm{mod.}\ p)\) ist oder nicht, und konnte ferner den Satz feststellen, dass für eine Primzahl \(p=4n+1\) stets ganze durch \(p\) nicht teilbare Zahlen \(x,\,y\) gefunden werden können, für welche \(x^2+y^2\) durch \(p\) teilbar wird\begingroup\renewcommand{\thefootnote}{**)}\footnote{Novi Comment. Acad. Petropol. \textbf{5} (1754/56) 1760, S. 3, Opera Omnia, Ser. I, vol. 2, S. 328.}\endgroup.

Auf festerem und breiterem Grunde baute \textsc{Gauss} diese Sätze auf. Dem schon Gesagten gemäss gibt es in der obigen Reihe der Potenzen von \(g\) eine solche Potenz \(g^\alpha\), dass \(a\equiv g^\alpha\ (\mathrm{mod.}\ p)\); der Exponent \(\alpha\) heisst nach \textsc{Gauss} der Index von \(a\). Aus den Gesetzen, welche diese Indizes befolgen, lässt sich leicht ein Kriterium ableiten, nach welchem die allgemeine binomische Kongruenz
\[
x^n\equiv a\ (\mathrm{mod.}\ p)
\]
auflösbar ist: versteht man unter \(\delta\) den grössten gemeinsamen Teiler von \(n\) und \(p-1\), so muss
\[
a^{\frac{p-1}{\delta}}\equiv1\ (\mathrm{mod.}\ p)
\]
sein. Für \(n=2\) ergibt sich so von selbst das \textsc{Euler}sche Kriterium. Gehört ferner eine Zahl \(a\) zum Exponenten \(d\), so bilden die allein einander inkongruenten Potenzen \(1,\,a,\,a^2,\ldots,a^{d-1}\) derselben die sogenannte Periode der Zahl \(a\ (\mathrm{mod.}\ p)\) und ihr Produkt ist \(\equiv-(-1)^d\). Wird für \(a\) eine primitive
\HistoricalFoot{}{2*}
\clearpage\SourceLeaf{19}\ArithmeticHead{12}{L}{BACHMANN, ÜBER GAUSS’ ZAHLENTHEORETISCHE ARBEITEN.}
\noindent Wurzel \(g\) gewählt, so findet sich hieraus sogleich der bereits von \textsc{Waring}\footnote{\textit{Meditationes algebraicae}, Cantabrigae 1770, S. 218.} mitgeteilte und \textsc{Wilson} zugeschriebene Satz, nach welchem das Produkt aller zu \(p\) primen Reste, nämlich
\[
1\,.\,2\,.\,3\,.\,.\,.(p-1)\equiv-1\ (\mathrm{mod.}\ p)
\]
ist. Diesen Satz bewies (unter Annahme der Existenz einer primitiven Wurzel) auf dem zuletzt angegebenen Wege vor \textsc{Gauss} schon \textsc{Euler}\begingroup\renewcommand{\thefootnote}{**)}\footnote{Opuscula analytica I (Petropoli 1783), S. 329.}\endgroup\ und aus wesentlich anderen Prinzipien \textsc{Lagrange}\begingroup\renewcommand{\thefootnote}{***)}\footnote{Nouv. Mémoires de l’Acad. Berlin (1771) 1773, \textbf{2}, S. 125, Oeuvres III, S. 425.}\endgroup. \textsc{Gauss} aber gab noch einen neuen Beweis des Satzes, der, auf die sogenannten bezüglich \(p\) assoziierten Zahlen gegründet, den Vorzug hat, dass er sich auf den allgemeineren Fall eines zusammengesetzten Moduls übertragen lässt. Der von \textsc{Gauss} nur ausgesprochene, hier geltende verallgemeinerte \textsc{Wilson}sche Satz wurde später auf dem angedeuteten Wege von \textsc{Brennecke}\begingroup\renewcommand{\thefootnote}{†)}\footnote{\textsc{Crelles} Journal für Mathematik \textbf{19} (1839), S. 319.}\endgroup, dann von \textsc{Schering}\begingroup\renewcommand{\thefootnote}{††)}\footnote{Acta Mathematica \textbf{1} (1883), S. 153.}\endgroup\ bewiesen.
